% Ελληνική Περίληψη

\begin{abstract}\setlanguage{greek}

Η επιλογή της κατάλληλης τοποθεσίας αποτελεί κρίσιμο παράγοντα επιτυχίας για κάθε επιχειρηματική δραστηριότητα. Σκοπός της παρούσας διπλωματικής μελέτης είναι η ανάπτυξη ενός συστήματος που συνδυάζει τη δύναμη των χωρικών δεδομένων με τις δυνατότητες της γενετικής τεχνητής νοημοσύνης, με στόχο την υποστήριξη της βέλτιστης επιλογής επιχειρηματικής τοποθεσίας. Το προτεινόμενο σύστημα αξιοποιεί γεωχωρικά χαρακτηριστικά, όπως πυκνότητα πληθυσμού, προσβασιμότητα, αλλά και δημογραφικά ή οικονομικά δεδομένα περιοχών. Μέσω ενός συνδυασμού χωρικής ανάλυσης και προηγμένων γλωσσικών μοντέλων (LLMs), το σύστημα αντλεί πληροφορίες από ετερογενείς πηγές, τις επεξεργάζεται και παράγει επεξηγηματικές και εξατομικευμένες προτάσεις τοποθεσιών για νέες επιχειρήσεις. Η μελέτη περιλαμβάνει την κατασκευή δυο pipelines βασισμένα σε LLMs, την επιλογή κατάλληλων γεωχωρικών χαρακτηριστικών, τη δημιουργία δείκτη αναζήτησης, την εκπαίδευση μοντέλου αξιολόγησης τοποθεσιών και την τελική παραγωγή απαντήσεων. Παράλληλα, αξιολογείται η αποτελεσματικότητα του συστήματος μέσω μετρικών ακρίβειας.

\vspace*{\fill}
\noindent{\large\bf{Λέξεις-κλειδιά:}}\\ 
όρος, όρος, $\dots$, όρος
\end{abstract}







